% Standalone problem document, generated from the adjacent README.md.
% Compile with XeLaTeX. Mathematical content is maintained in Markdown.
\documentclass[11pt,a4paper]{article}
\usepackage[fontsize=10.5pt]{scrextend}
\usepackage[margin=22mm,headheight=15pt,headsep=9mm,footskip=11mm]{geometry}
\usepackage{fontspec}
\setmainfont{texgyrepagella}[Extension=.otf,UprightFont=*-regular,BoldFont=*-bold,ItalicFont=*-italic,BoldItalicFont=*-bolditalic]
\setsansfont{texgyreheros}[Extension=.otf,UprightFont=*-regular,BoldFont=*-bold,ItalicFont=*-italic,BoldItalicFont=*-bolditalic]
\setmonofont{texgyrecursor}[Extension=.otf,UprightFont=*-regular,BoldFont=*-bold,ItalicFont=*-italic,BoldItalicFont=*-bolditalic,Scale=MatchLowercase]
\usepackage{amsmath,amssymb}
\usepackage{unicode-math}
\setmathfont{texgyrepagella-math.otf}
\usepackage{xcolor,microtype,enumitem,needspace,fancyhdr}
\usepackage{bookmark}
\definecolor{ink}{HTML}{17344B}
\definecolor{accent}{HTML}{197D80}
\definecolor{muted}{HTML}{596773}
\hypersetup{colorlinks=true,linkcolor=accent,urlcolor=accent,pdftitle={RA-04 - Clustered
singular-value gaps in randomized block Krylov
approximation},pdfauthor={Open Problems in Numerical Linear Algebra}}
\urlstyle{same}
\setlength{\parindent}{0pt}
\setlength{\parskip}{5pt plus 1pt minus 1pt}
\linespread{1.04}
\setlength{\emergencystretch}{3em}
\widowpenalty=10000
\clubpenalty=10000
\displaywidowpenalty=10000
\allowdisplaybreaks[1]
\setlist{nosep,leftmargin=1.5em}
\providecommand{\tightlist}{\setlength{\itemsep}{0pt}\setlength{\parskip}{0pt}}
\providecommand{\pandocbounded}[1]{#1}
\pagestyle{fancy}
\fancyhf{}
\fancyhead[L]{\sffamily\footnotesize\color{muted}OPEN PROBLEMS IN NUMERICAL LINEAR ALGEBRA}
\fancyhead[R]{\sffamily\bfseries\footnotesize\color{accent}RA-04}
\fancyfoot[L]{\parbox[b]{0.9\textwidth}{\sffamily\fontsize{8}{10}\selectfont\color{muted}Editorial ratings; a bounded search cannot certify that a problem remains unsolved.\\Literature check: 2026-09-10}}
\fancyfoot[R]{\sffamily\footnotesize\color{muted}\thepage}
\renewcommand{\headrulewidth}{0.3pt}
\renewcommand{\headrule}{\hbox to\headwidth{\color{accent}\leaders\hrule height \headrulewidth\hfill}}
\makeatletter
\renewcommand{\section}{\Needspace{5\baselineskip}\@startsection{section}{1}{0pt}{12pt plus 2pt minus 2pt}{4pt}{\sffamily\bfseries\large\color{ink}}}
\renewcommand{\subsection}{\Needspace{4\baselineskip}\@startsection{subsection}{2}{0pt}{9pt plus 1pt minus 1pt}{3pt}{\sffamily\bfseries\normalsize\color{ink}}}
\makeatother
\setcounter{secnumdepth}{0}
\begin{document}
{\sffamily\small\color{muted}Randomized and low-rank approximation\par}
\vspace{3pt}
{\raggedright\hyphenpenalty=10000\exhyphenpenalty=10000\sffamily\bfseries\fontsize{21}{24}\selectfont\color{ink}Clustered
singular-value gaps in randomized block Krylov approximation\par}
\vspace{7pt}
{\sffamily\small\textbf{Difficulty:} hard\quad\textbf{Importance:} interesting
to specialist\par}
{\sffamily\small\bfseries\color{accent}Status: Open\par}
\vspace{4pt}
{\color{accent}\hrule height 0.8pt}
\vspace{6pt}
\textbf{Rating rationale:} Hard because this is a focused spectral-gap
refinement of an established Krylov bound; specialist impact reflects
the source's assessment that its practical gain is limited.

Let \(A\in\mathbb R^{n\times d}\) have singular values
\(\sigma_1\geq\sigma_2\geq\cdots\). Fix integers \(1\leq b\leq k\), set
\(t=\lceil k/b\rceil\), and set \(k'=bt\leq\operatorname{rank}(A)\).
Define \[
\Delta_{k'}^{(b)}=\min_{1\leq i\leq k'-b}
\frac{\sigma_i^2-\sigma_{i+b}^2}{\sigma_i^2},
\] with the empty minimum set to one, and suppose this gap is positive.
Draw \(G\in\mathbb R^{n\times b}\) with independent standard Gaussian
entries. In exact arithmetic, let \(Z\) be an orthonormal basis of \[
\operatorname{range}[G,(AA^T)G,\ldots,(AA^T)^{q-1}G],
\qquad \widehat A=Z[Z^TA]_k,
\] where \([B]_k\) denotes a best rank-\(k\) approximation obtained by
truncating the SVD of \(B\).

Does an absolute constant \(C\) exist such that, for every such input
and every \(0<\varepsilon,\delta<1/2\), taking \[
q=\left\lceil C\left[
\frac{t}{\sqrt\varepsilon}\log\frac{2}{\Delta_{k'}^{(b)}}+
\frac1{\sqrt\varepsilon}\log\frac{n}{\delta\varepsilon}
\right]\right\rceil
\] gives, with probability at least \(1-\delta\), both \[
\|A-\widehat A\|_\xi\leq(1+\varepsilon)\|A-[A]_k\|_\xi
\quad(\xi=2,F),
\] and \(|\|A v_i\|_2^2-\sigma_i^2|\leq\varepsilon\sigma_{k+1}^2\) for
the ordered top \(k\) right singular vectors \(v_i\) of \(\widehat A\)?
Here \(\sigma_{k+1}=0\) if necessary. The factor two makes the logarithm
meaningful at a unit gap.

This is the gap-independent algorithmic component of the authors'
concluding conjecture. They also conjecture the corresponding
improvement in their gap-dependent theorem and random Krylov matrix
conditioning theorem; those are not separately counted here. The
displayed parameter regime is where the source's spectral-gap expression
is finite and defined.

The existing theorem instead depends logarithmically on every
consecutive gap and a leading spectral condition number. The proposed
bound would explain why a block can accommodate clusters of up to \(b\)
singular values. It does not ask for a finite-precision extension.

\newpage
\subsection{References}\label{references}

\begin{enumerate}
\def\labelenumi{\arabic{enumi}.}
\tightlist
\item
  T. Chen, E. N. Epperly, R. A. Meyer, C. Musco, and A. Rao, \emph{Does
  block size matter in randomized block Krylov low-rank approximation?},
  SODA 2026, pp.~1026--1046
  (\href{https://doi.org/10.1137/1.9781611978971.42}{published paper});
  arXiv:2508.06486v2. Section 5 states the conjecture; Section 3.3 gives
  the full quantitative Theorem 1.3, and Algorithm 1 fixes the
  algorithm. \href{https://arxiv.org/html/2508.06486v2}{Paper}.
\item
  R. A. Meyer, C. Musco, and C. Musco, \emph{On the Unreasonable
  Effectiveness of Single Vector Krylov Methods for Low-Rank
  Approximation}, SODA 2024, pp.~811--845, especially Theorem 4.5.
  \href{https://arxiv.org/abs/2305.02535}{Paper}.
\end{enumerate}

\subsection{Status check --- 2026-09-08}\label{status-check-2026-09-08}

Searched the exact 2025 paper title, ``block Krylov b-th order gap
conjecture'', and ``randomized block Krylov clustered gaps 2026'';
checked the current arXiv abstract/version record and Section 5. No
proof or counterexample was found. Input perturbation results in Section
3.5 do not establish the unperturbed statement above.

\subsection{Audit --- 2026-09-10}\label{audit-2026-09-10}

Rechecked \href{https://arxiv.org/html/2508.06486v2}{§5 of the latest
arXiv v2}, which retains the clustered-gap conjecture, and added the
\href{https://doi.org/10.1137/1.9781611978971.42}{SODA 2026
publication}, pp.~1026--1046. Clustered-gap and later low-rank-query
searches found no resolution. Randomly perturbing the input changes the
displayed target.
\end{document}
